% A content fragment, meant to be \input{} or \include{}d from a master
% document that supplies \documentclass, packages (amsmath, unicode-math +
% fontspec for the raw Unicode math symbols below — compile the master with
% xelatex or lualatex), \begin{document}/\end{document}, and its own
% \bibliographystyle/\bibliography pointing at "9 References.bib".

\section{Dietary Requirements}

The remaining four bands — protein, dietary fiber, fat, and carbohydrate — set healthy daily ranges for diet composition rather than for the size of the energy deficit itself, each scaled to lean mass or to intake, in the four subsections below. A final subsection then bounds the deficit itself, both in rate and in magnitude relative to BMR.

\subsection{Protein Requirements}

Protein needs for energy-restricted, resistance-trained individuals with low body fat are estimated at a floor p\_min and ceiling p\_max of 2.3–3.1 g per kg of fat-free (lean) mass per day, scaled upward with the severity of caloric restriction and leanness \cite{ref_helms2014}. The corresponding lower and upper healthy daily protein amounts, P\_min and P\_max, are given in eqn.~\eqref{eqn_pmin} and eqn.~\eqref{eqn_pmax}.

\begin{equation}
\label{eqn_pmin}
P_{min} = p_{min} × LBM
\end{equation}
\begin{equation}
\label{eqn_pmax}
P_{max} = p_{max} × LBM
\end{equation}

\subsection{Dietary Fiber Requirements}

Dietary fiber intake is recommended relative to energy intake via a floor coefficient b\_min of roughly 14 g per 1,000 kcal consumed, while a practical upper ceiling b\_max is instead scaled to body weight, per the USDA Dietary Guidelines for Americans \cite{ref_usda2020}. These two different bases give the lower and upper healthy daily dietary fiber amounts, F\_min and F\_max, in eqn.~\eqref{eqn_fmin} and eqn.~\eqref{eqn_fmax} respectively.

\begin{equation}
\label{eqn_fmin}
F_{min} = b_{min} × \frac{TEI}{1000}
\end{equation}
\begin{equation}
\label{eqn_fmax}
F_{max} = b_{max} × m̄
\end{equation}

\subsection{Fat Requirements}

Fat intake is recommended within an Acceptable Macronutrient Distribution Range (AMDR) of 20–35\% of total energy intake for adults \cite{ref_iom2005}, the floor and ceiling percentages denoted k\_min and k\_max. Converted to grams at fat's fixed energy density of 9 kcal/g, this range gives the lower and upper healthy daily fat amounts, G\_min and G\_max, in eqn.~\eqref{eqn_gmin} and eqn.~\eqref{eqn_gmax}.

\begin{equation}
\label{eqn_gmin}
G_{min} = \frac{(k_{min}/100) × TEI}{9}
\end{equation}
\begin{equation}
\label{eqn_gmax}
G_{max} = \frac{(k_{max}/100) × TEI}{9}
\end{equation}

\subsection{Carbohydrate Requirements}

Carbohydrate intake carries the same AMDR structure, recommended at 45–65\% of total energy intake for adults \cite{ref_iom2005}, the floor and ceiling percentages denoted q\_min and q\_max. Converted to grams at carbohydrate's fixed energy density of 4 kcal/g, this range gives the lower and upper healthy daily carbohydrate amounts, C\_min and C\_max, in eqn.~\eqref{eqn_cmin} and eqn.~\eqref{eqn_cmax}.

\begin{equation}
\label{eqn_cmin}
C_{min} = \frac{(q_{min}/100) × TEI}{4}
\end{equation}
\begin{equation}
\label{eqn_cmax}
C_{max} = \frac{(q_{max}/100) × TEI}{4}
\end{equation}

\subsection{Deficit Safety Limits}

Beyond composition, the pace and depth of the deficit itself carry risk: too fast a rate of fat loss, or too deep a cut relative to BMR, accelerates lean-mass loss and metabolic adaptation. The weekly rate Δm\% (eqn.~\eqref{eqn_dmpc}) should stay within 0.5–1\%, the range shown to preserve fat-free mass and strength/power performance in trained individuals under caloric restriction \cite{ref_garthe2011}. The daily deficit expressed as a share of BMR, D\_bmr\%, given in eqn.~\eqref{eqn_dbmr}, should stay under 20\%, in line with the moderate, non-aggressive deficits recommended to protect fat-free mass and resting energy expenditure during contest-style caloric restriction \cite{ref_helms2014prep}.

\begin{equation}
\label{eqn_dbmr}
D_{bmr}\% = 100 × \frac{BMR - TEI}{BMR}
\end{equation}
